\documentclass[12pt]{article}

\usepackage[utf8]{inputenc}
\usepackage[T1]{fontenc}
\usepackage[brazil]{babel}
\usepackage{geometry}
\geometry{a4paper, top=3.5cm, bottom=2.5cm, left=3.0cm, right=3.0cm}

\usepackage{times} % Fonte padrão SBC
\usepackage{graphicx}
\usepackage{amsmath, amssymb, amsfonts}
\usepackage{booktabs} % Tabelas profissionais
\usepackage{multirow}
\usepackage{float}
\usepackage{subcaption}
\usepackage{xcolor}
\usepackage{listings}
\usepackage{hyperref}

% Configuração de links
\hypersetup{
    colorlinks=true,
    linkcolor=black,
    citecolor=black,
    urlcolor=blue
}

% Configuração para exibição de trechos de código (se necessário)
\lstset{
    basicstyle=\ttfamily\footnotesize,
    breaklines=true,
    frame=single,
    numbers=left,
    numberstyle=\tiny\color{gray},
    keywordstyle=\color{blue}\bfseries,
    commentstyle=\color{teal},
    stringstyle=\color{orange}
}

% Estilização de seções no padrão SBC
\usepackage{titlesec}
\titleformat{\section}{\normalfont\large\bfseries}{\thesection.}{0.5em}{}
\titleformat{\subsection}{\normalfont\normalsize\bfseries}{\thesubsection.}{0.5em}{}

% ==============================================================================
% METADADOS DO ARTIGO
% ==============================================================================
\title{\textbf{Análise Empírica, Auditoria Algorítmica e Otimização para o Problema do Caixeiro Viajante Geométrico}}

\author{
    Rafael Junqueira Baraacat Silva, Cauã Neto, Renato Villela \\
    \small Departamento de Ciência da Computação -- Universidade Federal de Minas Gerais (UFMG) \\
    \small Belo Horizonte -- MG -- Brasil \\
    \small \texttt{junqueira.rafael149@gmail.com}
    \small \texttt{email caua}
    \small \texttt{email renato}
}


\begin{document}

\maketitle

\noindent\textbf{Resumo.} Este relatorio apresenta uma análise empírica, auditoria algorítmica e otimização de desempenho aplicadas ao Problema do Caixeiro Viajante (TSP) Geométrico em instâncias euclidianas bidimensionais da biblioteca TSPLIB. Avaliam-se um algoritmo exato baseado em \textit{Branch-and-Bound} e dois algoritmos aproximados clássicos: \textit{Twice-around-the-tree} e \textit{Christofides}. Adicionalmente, investiga-se o impacto na poda de cinco estratégias de inicialização de limite superior (\textit{warm-start}): \textit{Cold Start}, Vizinho Mais Próximo, Inserção por Fecho Convexo, \textit{Twice-around-the-tree} e \textit{Christofides}. Por meio de \textit{profiling} em nível de sub-rotina, modelagem estatística por regressão assintótica e análise de sensibilidade topológica, caracterizam-se gargalos computacionais e identificam-se atributos geométricos que regem o \textit{gap} de aproximação e o colapso da busca exata.

% ==============================================================================
% 1. INTRODUÇÃO
% ==============================================================================
\section{Introdução}

% TODO: Contextualizar o Problema do Caixeiro Viajante (TSP) Geométrico (EUC_2D),
% sua classificação na classe NP-Difícil e a relevância do compromisso entre
% otimalidade garantida (métodos exatos) e escalabilidade (heurísticas e aproximações).

% TODO: Apresentar os objetivos específicos do trabalho prático (DCC207 - Algoritmos 2):
% 1) Auditoria, refatoração e instrumentação de código;
% 2) Avaliação do impacto de 5 estratégias de warm-start na poda do Branch-and-Bound;
% 3) Modelagem empírica de complexidade temporal e espacial via regressão;
% 4) Análise de sensibilidade topológica e identificação de instâncias adversariais.

% TODO: Descrever brevemente a organização das demais seções do artigo.

% ==============================================================================
% 2. IMPLEMENTAÇÃO, PROFILING E AUDITORIA
% ==============================================================================
\section{Implementação, Profiling e Auditoria}

\subsection{Arquitetura e Estruturas de Dados Adotadas}
% TODO: Descrever a linguagem utilizada (Python3 ou C++11), representação do grafo
% (matriz de adjacência/distâncias pré-computada vs. cálculo sob demanda) e
% estruturas auxiliares (filas de prioridade/heaps, conjuntos disjuntos/Union-Find,
% listas de adjacência para grafos eulerianos, etc.).

\subsection{Estratégias de Busca e Limite Inferior no Branch-and-Bound}
% TODO: Detalhar a estratégia de exploração da árvore de busca (Best-First Search,
% Depth-First Search ou híbrida) e a função de limite inferior (lower bound)
% utilizada (ex: soma das duas menores arestas incidentes / 1-tree / MST).
% Explicar como o estado parcial do caminho é armazenado sem cópias excessivas de memória.

\subsection{Implementação das Heurísticas e Algoritmos Aproximados}
% TODO: Explicar brevemente as 5 estratégias de inicialização (warm-start):
% (1) Cold Start / Trivial (1 -> 2 -> ... -> n -> 1)
% (2) Vizinho Mais Próximo (Nearest Neighbor)
% (3) Inserção Baseada em Fecho Convexo (Convex Hull Insertion)
% (4) Twice-around-the-tree (MST + caminhamento em pré-ordem)
% (5) Christofides (MST + emparelhamento perfeito de peso mínimo nos vértices de grau ímpar + circuito euleriano + shortcutting)

\subsection{Auditoria de Código, Refatoração e Gargalos de Sub-rotina}
% TODO: Relatar o processo de auditoria e refatoração (incluindo eventuais
% otimizações sobre código base ou gerado por LLMs):
% - Eliminação de alocações redundantes de matrizes e impacto de garbage collection/memória;
% - Resultados do profiling em nível de sub-rotina: isolar o custo computacional de
%   etapas críticas como a construção da Árvore Geradora Mínima (MST) e o
%   Emparelhamento Perfeito de Peso Mínimo (Minimum Weight Perfect Matching) no Christofides.

% ==============================================================================
% 3. METODOLOGIA EXPERIMENTAL
% ==============================================================================
\section{Metodologia Experimental}

\subsection{Ambiente Computacional e Conjunto de Teste (TSPLIB)}
% TODO: Especificar hardware (CPU, RAM, SO) e versões de compilador/interpretador e bibliotecas.
% Descrever a seleção das 12 a 16 instâncias EUC_2D da TSPLIB (N entre 10 e 75):
% - Faixa Principal (N <= 40, ~80% das instâncias, ex: ulysses16, gr24, fri26, bays29)
% - Faixa Limite / Estresse (N in [45, 75], ex: swiss42, att48, etc.)
% Mencionar o timeout de 30 minutos (1800s) e o tratamento de execuções abortadas (NA).

\subsection{Extração de Features Geométricas e Topológicas}
% TODO: Definir formalmente as métricas extraídas de cada instância para a Análise de Sensibilidade:
% - Razão de aspecto (aspect ratio) do bounding box;
% - Proporção de vértices pertencentes ao fecho convexo;
% - Entropia da distribuição espacial ou densidade de clustering;
% - Estatísticas da MST (ex: grau máximo de nó / presença de hubs);
% - Coeficiente de variação (CV) das distâncias ao vizinho mais próximo.

% ==============================================================================
% 4. EXPERIMENTOS E RESULTADOS (MODELAGEM)
% ==============================================================================
\section{Experimentos e Resultados}

\subsection{Visão Geral de Desempenho e Qualidade das Soluções}
% TODO: Preencher com os dados consolidados das execuções (tempos, custos, gaps e NAs).

\begin{table}[H]
\centering
\caption{Resumo comparativo de tempo de execução (s) e \textit{gap} de otimalidade (\%) nas instâncias selecionadas da TSPLIB (limite de 1800s para métodos exatos).}
\label{tab:resultados_gerais}
\resizebox{\textwidth}{!}{%
\begin{tabular}{@{}lcccccccc@{}}
\toprule
\multirow{2}{*}{\textbf{Instância}} & \multirow{2}{*}{\textbf{$N$}} & \multirow{2}{*}{\textbf{Ótimo}} & \multicolumn{2}{c}{\textbf{Branch-and-Bound (Melhor WS)}} & \multicolumn{2}{c}{\textbf{Twice-around-the-tree}} & \multicolumn{2}{c}{\textbf{Christofides}} \\ \cmidrule(lr){4-5} \cmidrule(lr){6-7} \cmidrule(lr){8-9}
 &  &  & \textbf{Tempo (s)} & \textbf{Nós Exp.} & \textbf{Tempo (s)} & \textbf{Gap (\%)} & \textbf{Tempo (s)} & \textbf{Gap (\%)} \\ \midrule
\texttt{ulysses16} & 16 & -- & -- & -- & -- & -- & -- & -- \\
\texttt{gr24}      & 24 & -- & -- & -- & -- & -- & -- & -- \\
\texttt{fri26}     & 26 & -- & -- & -- & -- & -- & -- & -- \\
\texttt{bays29}    & 29 & -- & -- & -- & -- & -- & -- & -- \\
\texttt{swiss42}   & 42 & -- & -- & -- & -- & -- & -- & -- \\
\texttt{att48}     & 48 & -- & -- & -- & -- & -- & -- & -- \\
\bottomrule
\end{tabular}%
}
\end{table}

\subsection{Profiling de Sub-rotinas e Modelagem de Complexidade Empírica}
% TODO: Apresentar os modelos de regressão ajustados (SciPy / Scikit-Learn) para tempo
% e memória em função de N (ex: ajustes polinomiais O(N^2), O(N^3) para Christofides
% e exponenciais/fatoriais para B&B), relatando R^2 e RMSE.
% Discutir discrepâncias entre complexidade teórica e empírica (overhead de memória/cache).

\begin{table}[H]
\centering
\caption{Parâmetros das curvas de regressão ajustadas para complexidade temporal e espacial.}
\label{tab:regressao_complexidade}
\begin{tabular}{@{}lcccc@{}}
\toprule
\textbf{Algoritmo / Sub-rotina} & \textbf{Modelo Teórico} & \textbf{Curva Ajustada} & \textbf{$R^2$} & \textbf{RMSE} \\ \midrule
Branch-and-Bound & $O(N!)$ / $O(b^N)$ & -- & -- & -- \\
Twice-around-the-tree & $O(N^2 \log N)$ & -- & -- & -- \\
Christofides (Total) & $O(N^3)$ & -- & -- & -- \\
\quad -- Sub-rotina MST & $O(N^2)$ & -- & -- & -- \\
\quad -- Sub-rotina Matching & $O(N^3)$ & -- & -- & -- \\
\bottomrule
\end{tabular}
\end{table}

% TODO: Inserir figuras das curvas de regressão e breakdown de tempo por sub-rotina.
% \begin{figure}[H]
%     \centering
%     \includegraphics[width=0.85\textwidth]{figuras/regressao_tempo.pdf}
%     \caption{Ajuste empírico do tempo de execução em função do número de vértices $N$.}
%     \label{fig:regressao_tempo}
% \end{figure}

\subsection{Impacto do Warm-Start na Eficiência de Poda}
% TODO: Comparar nós expandidos, nós podados e tempo total das 5 estratégias de warm-start.
% Analisar o trade-off: o tempo investido na heurística inicial (ex: Christofides vs. Nearest Neighbor)
% compensa a redução da árvore de busca?

\begin{table}[H]
\centering
\caption{Comparação do impacto das estratégias de \textit{warm-start} na árvore de busca do Branch-and-Bound.}
\label{tab:warm_start}
\resizebox{\textwidth}{!}{%
\begin{tabular}{@{}lccccc@{}}
\toprule
\textbf{Estratégia de Inicialização} & \textbf{Gap Inicial (\%)} & \textbf{Tempo Heur. (s)} & \textbf{Nós Expandidos} & \textbf{Nós Podados} & \textbf{Tempo Total B\&B (s)} \\ \midrule
1. Cold Start / Trivial & -- & -- & -- & -- & -- \\
2. Nearest Neighbor & -- & -- & -- & -- & -- \\
3. Convex Hull Insertion & -- & -- & -- & -- & -- \\
4. Twice-around-the-tree & -- & -- & -- & -- & -- \\
5. Christofides & -- & -- & -- & -- & -- \\
\bottomrule
\end{tabular}%
}
\end{table}

\subsection{Análise de Sensibilidade Topológica e Instâncias Adversariais}
% TODO: Apresentar o modelo estatístico (regressão múltipla, árvore de decisão ou
% importância de features) correlacionando as características geométricas da instância
% (aspect ratio, % no fecho convexo, entropia, grau máximo da MST, CV das distâncias)
% com o gap de aproximação e a dificuldade para o Branch-and-Bound.
% Discutir casos adversariais específicos encontrados no conjunto de teste.

% \begin{figure}[H]
%     \centering
%     \includegraphics[width=0.85\textwidth]{figuras/analise_sensibilidade.pdf}
%     \caption{Importância das features topológicas e correlação com o gap de aproximação.}
%     \label{fig:sensibilidade}
% \end{figure}

% ==============================================================================
% 5. CONCLUSÕES
% ==============================================================================
\section{Conclusões}

% TODO: Sintetizar as principais descobertas empíricas:
% - Ponto de colapso observado para o Branch-and-Bound (Faixa Principal vs. Faixa Limite);
% - Melhor custo-benefício entre as estratégias de warm-start;
% - Gargalos dominantes identificados no profiling (ex: Matching vs. MST);
% - Como a geometria/topologia dos pontos afeta o desempenho das heurísticas e aproximações.

% ==============================================================================
% REFERÊNCIAS BIBLIOGRÁFICAS
% ==============================================================================
\begin{thebibliography}{9}

\bibitem{christofides1976}
Christofides, N. (1976).
\newblock Worst-case analysis of a new heuristic for the travelling salesman problem.
\newblock \emph{Report 388, Graduate School of Industrial Administration, CMU}.

\bibitem{reinelt1991}
Reinelt, G. (1991).
\newblock TSPLIB -- A Traveling Salesman Problem Library.
\newblock \emph{ORSA Journal on Computing}, 3(4):376--384.

\bibitem{cormen2009}
Cormen, T. H., Leiserson, C. E., Rivest, R. L., and Stein, C. (2009).
\newblock \emph{Introduction to Algorithms}.
\newblock MIT Press, 3rd edition.

\end{thebibliography}

% ==============================================================================
% APÊNDICE OBRIGATÓRIO (CASO UTILIZE LLMs)
% ==============================================================================
\appendix
\section{Histórico de Prompts Utilizados (Uso de LLMs)}
% TODO: Conforme exigido na seção "O que entregar?" do enunciado, registrar aqui
% os prompts utilizados para geração/refatoração de código ou auxílio no trabalho,
% caso ferramentas de IA generativa tenham sido empregadas.

\end{document}
